
\documentclass[fontset=none,11pt,a4paper]{ctexart}
\usepackage[margin=1.8cm]{geometry}
\usepackage{fontspec}
\usepackage{amssymb}
\usepackage{enumitem}
\usepackage{hyperref}
\usepackage{xcolor}
\usepackage{titlesec}
\usepackage{parskip}
\IfFontExistsTF{Songti SC}
  {\setCJKmainfont{Songti SC}\setCJKsansfont{Songti SC}\setCJKmonofont{Songti SC}}
  {\setCJKmainfont[BoldFont=FandolSong-Bold.otf]{FandolSong-Regular.otf}\setCJKsansfont[BoldFont=FandolSong-Bold.otf]{FandolSong-Regular.otf}\setCJKmonofont[BoldFont=FandolSong-Bold.otf]{FandolSong-Regular.otf}}
\IfFontExistsTF{TeX Gyre Termes}{\setmainfont{TeX Gyre Termes}}{\setmainfont{Times New Roman}}
\hypersetup{colorlinks=true, linkcolor=black, urlcolor=blue}
\setlength{\parindent}{0pt}
\setlist[itemize]{leftmargin=2em,itemsep=0.15em}
\setlist[enumerate]{leftmargin=2.6em,itemsep=0.15em}
\titleformat{\section}{\Large\bfseries}{\thesection.}{0.6em}{}
\titleformat{\subsection}{\large\bfseries}{\thesubsection}{0.6em}{}
\newcommand{\answerline}[1]{\underline{\hspace{#1}}}
\newcommand{\choicebox}[1]{$\square$~#1}
\newcommand{\blocktime}{\vspace{0.25em}\noindent 用时：\answerline{4cm}\par\vspace{0.6em}}
\newcommand{\tightpagebreak}{\vspace{0.4em}\hrule\vspace{0.6em}}

\title{\textbf{AIH 人机实验二：子 Agent 诊断性评测（T1）}}
\author{}
\date{}

\begin{document}
\maketitle
\vspace{-2em}

\textbf{作答说明：}本表单包含三部分：Block A 句子级诊断、Block B TimeBlock 时间推理、Block C 跨文本对齐。请只根据题面材料作答；不需要查阅整篇原文。Block A 的“时间信息原文”请填写原句中直接表达时间的短语，若没有则填“无”。

\section{Block A：句子级诊断（Agent 2--4）}
\blocktime
覆盖 Agent 2 时间信息原文抽取、Agent 3 下沉句判断、Agent 4 插叙/倒叙判断。

\paragraph{A-1-01 \quad 【7.1.2】}
\textbf{上下文（前五句 / 目标句 / 后五句）：}

\noindent 【7.1.1】项籍是下相人，字羽。\par
\noindent \textcolor{blue}{\textbf{【7.1.2】开始起兵时二十四岁。}}\par
\noindent 【7.1.3】他的叔父是项梁，项梁的父亲就是楚将项燕，被秦将王翦所杀的那个人。\par
\noindent 【7.1.4】项氏世代为楚将，封于项，所以姓项氏。\par
\noindent 【7.2.1】项籍小时候，学习认字写字，没有学成。\par
\noindent 【7.2.2】放弃了学字，改学击剑，又没有学成。\par
\noindent 【7.2.3】项梁很生他的气。\par

\begin{itemize}
  \item 时间信息原文：\answerline{7.5cm}（无则填“无”）
  \item 是否为下沉句：\choicebox{是}\quad \choicebox{否}
  \item 是否为插叙/倒叙：\choicebox{是}\quad \choicebox{否}
\end{itemize}

\paragraph{A-1-02 \quad 【7.48.3】}
\textbf{上下文（前五句 / 目标句 / 后五句）：}

\noindent 【7.47.14】最后，郎中骑杨喜，骑司马吕马童，郎中吕胜、杨武各自得到了项王的一段肢体。\par
\noindent 【7.47.15】五个人把肢体合拢起来，都确实是项王的。\par
\noindent 【7.47.16】所以把准备封赏的土地分为五部分：封吕马童为中水侯，封王翳为杜衍侯，封杨喜为赤泉侯，封杨武为吴防侯，封吕胜为涅阳侯。\par
\noindent 【7.48.1】项王死后，楚国各地都投降了汉军，只有鲁城不肯投降。\par
\noindent 【7.48.2】汉王就带领天下士卒打算屠毁鲁城，因为他们坚守礼义，为主人以死守节，就拿项王的头给鲁城人看，鲁城父兄才投降了。\par
\noindent \textcolor{blue}{\textbf{【7.48.3】最初，楚怀王曾封项籍为鲁公，等到项籍死了，鲁城又最后投降，所以用鲁公的礼仪把项王理葬在谷城。}}\par
\noindent 【7.48.4】汉王为项王举哀，哭了一场，然后离开了鲁城。\par
\noindent 【7.49.1】各支项氏宗族，汉王都不诛杀。\par
\noindent 【7.49.2】封项伯为射阳侯。\par
\noindent 【7.49.3】桃侯、平景侯、玄武侯都是项氏宗族，赐姓刘。\par
\noindent 【7.50.1】太史公说：“我听周生说“舜的眼睛大概是两个瞳孔”，又听说项羽也是两个瞳孔。项羽难道是舜的后裔吗？为什么兴起得这么迅速啊！秦朝政治腐败，陈涉首先发难，豪杰蜂起，相互争夺，不可胜数。然而项羽毫无凭借，乘势起于民间，三年时间，就率领五路诸侯军消灭了秦朝，分割天下，封王建侯，政自己出，号为“霸王”，虽然没有始终保持他的地位，但近古以来，还未曾有过这样的事情。等到项羽放弃关中，怀恋楚地，放逐义帝而自立为王，抱怨王侯背叛自己，这时已经难以控制局势了。自我夸耀功勋，逞一己私智，不效法古人，以为创立霸王的事业，需要用武力来经营天下，终于五年时间覆灭了他自己的国家，身死东城，还没有觉悟，不自我谴责，这就不对了。竟然用“上天灭亡我，不是我用兵打仗的过错”为借口，难道不是太荒谬了吗！”\par

\begin{itemize}
  \item 时间信息原文：\answerline{7.5cm}（无则填“无”）
  \item 是否为下沉句：\choicebox{是}\quad \choicebox{否}
  \item 是否为插叙/倒叙：\choicebox{是}\quad \choicebox{否}
\end{itemize}

\paragraph{A-1-03 \quad 【8.5.11】}
\textbf{上下文（前五句 / 目标句 / 后五句）：}

\noindent 【8.5.6】看了鲁元，也是贵相。\par
\noindent 【8.5.7】老人已经走了，高祖正好从别人家来到田间，吕后告诉他一位客人从这里经过，给我们母子看相，说将来都是大贵人，高祖问老父在哪儿.吕后说：“走出不远。”\par
\noindent 【8.5.8】高祖追上了老人，向他询问。\par
\noindent 【8.5.9】老人说：“刚才相过夫人和孩子，他们都跟你相似，你的相貌，贵不可言。”\par
\noindent 【8.5.10】高祖便道谢说：“如果真像老父所说，决不忘记对我的恩德。”\par
\noindent \textcolor{blue}{\textbf{【8.5.11】等到高祖显贵，竟然不知道老人的去处了。}}\par
\noindent 【8.6.1】高祖做亭长，以竹皮为帽，这帽子是他派求盗到薛县制做的，经常戴着它。\par
\noindent 【8.6.2】等到显贵时，仍然常常戴着，所谓“刘氏冠”，就是指这种帽子。\par
\noindent 【8.7.1】高祖因身任亭长，为县里送役徒去郦山，役徒多在途中逃亡。\par
\noindent 【8.7.2】他估计，等走到郦山，大概都逃光了。\par
\noindent 【8.7.3】到丰邑西面的沼泽地带，停下来喝酒，夜间高祖就释放了所押送的役徒。\par

\begin{itemize}
  \item 时间信息原文：\answerline{7.5cm}（无则填“无”）
  \item 是否为下沉句：\choicebox{是}\quad \choicebox{否}
  \item 是否为插叙/倒叙：\choicebox{是}\quad \choicebox{否}
\end{itemize}

\paragraph{A-1-04 \quad 【8.6.2】}
\textbf{上下文（前五句 / 目标句 / 后五句）：}

\noindent 【8.5.8】高祖追上了老人，向他询问。\par
\noindent 【8.5.9】老人说：“刚才相过夫人和孩子，他们都跟你相似，你的相貌，贵不可言。”\par
\noindent 【8.5.10】高祖便道谢说：“如果真像老父所说，决不忘记对我的恩德。”\par
\noindent 【8.5.11】等到高祖显贵，竟然不知道老人的去处了。\par
\noindent 【8.6.1】高祖做亭长，以竹皮为帽，这帽子是他派求盗到薛县制做的，经常戴着它。\par
\noindent \textcolor{blue}{\textbf{【8.6.2】等到显贵时，仍然常常戴着，所谓“刘氏冠”，就是指这种帽子。}}\par
\noindent 【8.7.1】高祖因身任亭长，为县里送役徒去郦山，役徒多在途中逃亡。\par
\noindent 【8.7.2】他估计，等走到郦山，大概都逃光了。\par
\noindent 【8.7.3】到丰邑西面的沼泽地带，停下来喝酒，夜间高祖就释放了所押送的役徒。\par
\noindent 【8.7.4】高祖说：“各位都走吧，我也从此一去不返了！”\par
\noindent 【8.7.5】役徒中有十多个年轻力壮的愿意跟随高祖。\par

\begin{itemize}
  \item 时间信息原文：\answerline{7.5cm}（无则填“无”）
  \item 是否为下沉句：\choicebox{是}\quad \choicebox{否}
  \item 是否为插叙/倒叙：\choicebox{是}\quad \choicebox{否}
\end{itemize}

\paragraph{A-1-05 \quad 【8.9.1】}
\textbf{上下文（前五句 / 目标句 / 后五句）：}

\noindent 【8.8.2】高祖怀疑这件事与自己有关，就逃跑藏了起来，隐身在芒山、砀山一带的山泽岩石之间。\par
\noindent 【8.8.3】吕后和别人一块儿寻找，常常一去就找到了高祖。\par
\noindent 【8.8.4】高祖感到奇怪，就问吕后，吕后说：“你所处的地方上面常有云气，向着有云气的地方去找，常常可以找到你。”\par
\noindent 【8.8.5】高祖心里非常高兴。\par
\noindent 【8.8.6】沛县子弟有的听到这件事.很多人都想归附他。\par
\noindent \textcolor{blue}{\textbf{【8.9.1】秦二世元年秋天，陈胜等在蕲县起义，到了陈县自立为王，号称“张楚”。}}\par
\noindent 【8.9.2】各郡县都大多杀死长官，响应陈胜。\par
\noindent 【8.9.3】沛县县令恐惧，想要以沛具响应陈胜，主吏萧何、狱掾曹参对他说：“您身为秦朝的官吏，如今要叛秦起事，率领沛县子弟，恐怕他们不愿听命。希望您召集逃亡在外面的人,可以得到几百人。利用这股力量胁持群众，群众不敢不听您的命令。”\par
\noindent 【8.9.4】县令就派樊哙去召唤刘季，刘季的队伍已经近百人了。\par
\noindent 【8.10.1】于是樊哙跟着刘季来到沛县。\par
\noindent 【8.10.2】沛县县令又后悔了，恐怕刘季发生变故，就关闭城门，派人防守，打算杀掉萧何、曹参。\par

\begin{itemize}
  \item 时间信息原文：\answerline{7.5cm}（无则填“无”）
  \item 是否为下沉句：\choicebox{是}\quad \choicebox{否}
  \item 是否为插叙/倒叙：\choicebox{是}\quad \choicebox{否}
\end{itemize}

\paragraph{A-1-06 \quad 【8.10.7】}
\textbf{上下文（前五句 / 目标句 / 后五句）：}

\noindent 【8.10.2】沛县县令又后悔了，恐怕刘季发生变故，就关闭城门，派人防守，打算杀掉萧何、曹参。\par
\noindent 【8.10.3】萧何、曹参恐惧，翻过城墙依附刘季。\par
\noindent 【8.10.4】刘季用帛写了一封信，射到城上，告诉沛县父老说：“天下苦于秦朝的暴政已经很久了。现在父老为沛令守城，但各国诸侯都已起事，就要屠戮沛县。如果沛县父老共同起来杀死沛令，选择子弟中可以立为首领的做领导，以响应诸侯军，那就能保全身家性命。不然的话，父子全遭杀害，死得毫无意义。”\par
\noindent 【8.10.5】父老们就率领子弟共同杀了沛令，打开城门，迎接刘季，想让他做柿县县令。\par
\noindent 【8.10.6】刘季说：“天下正在混乱当中，诸侯都已起事，如果推选的将领不胜任，就会一败涂地。我不是吝惜自己的生命，只怕才劣力薄，不能保全父兄子弟。这是件大事，希望另外共同推选一位能够胜任的人。”\par
\noindent \textcolor{blue}{\textbf{【8.10.7】萧何、曹参等都是文官，看重身家性命，怕事情不成，秦朝会诛灭他们的全族，所以都推让刘季。}}\par
\noindent 【8.10.8】父老们都说：“我们平时听到刘季许多奇异的事情，看来刘季是该显贵的。而且又经过占卜，没有比刘季更吉利的。”\par
\noindent 【8.10.9】这时刘季再三谦让，大家都不敢担任，最后还是立刘季为沛公。\par
\noindent 【8.10.10】在沛县衙门的庭院里祭祀黄帝和蚩尤，又用牲血衅鼓旗。\par
\noindent 【8.10.11】旗子一律红色，因为刘季所杀蛇是白帝的儿子，杀蛇的是赤帝的儿子，所以崇尚赤色。\par
\noindent 【8.10.12】于是少年子弟和有势的官吏，如萧何、曹参、樊哙等人，都为沛公征集兵员，集合了两三千人，攻打胡陵、方与，回军固守丰邑。\par

\begin{itemize}
  \item 时间信息原文：\answerline{7.5cm}（无则填“无”）
  \item 是否为下沉句：\choicebox{是}\quad \choicebox{否}
  \item 是否为插叙/倒叙：\choicebox{是}\quad \choicebox{否}
\end{itemize}

\paragraph{A-1-07 \quad 【7.17.7】}
\textbf{上下文（前五句 / 目标句 / 后五句）：}

\noindent 【7.17.2】函谷关有兵把守，不能进去。\par
\noindent 【7.17.3】又听说沛公已经攻破咸阳，项羽大怒，派当阳君等扣关。\par
\noindent 【7.17.4】项羽便进入了函谷关，到达戏水西岸。\par
\noindent 【7.17.5】沛公驻军霸上，没有能够和项羽相见。\par
\noindent 【7.17.6】沛公左司马曹无伤派人对项羽说：“沛公想称王关中，使子婴为相，占有了全部珍宝。”\par
\noindent \textcolor{blue}{\textbf{【7.17.7】项羽怒气冲天他说：“明天早晨饱餐士卒，将击溃沛公的军队！”}}\par
\noindent 【7.17.8】这时，项羽有兵四十万，驻扎在新丰鸿门，沛公有兵十万，驻扎在霸上。\par
\noindent 【7.17.9】范增劝告项羽说：“沛公在山东时，贪财好货，喜爱美女。现在进了关，不收财物，不亲近妇女，由此看来，他的志向不小。我叫人观望他上空的云气，都呈龙虎形状，五颜六色，这是天子之气。赶快进击，不要失掉机会。”\par
\noindent 【7.18.1】楚国左尹项伯这个人，是项羽的叔父，一向和留侯张良相友好。\par
\noindent 【7.18.2】张良这时跟随着沛公，项伯就夜间骑马跑到沛公军营，私下见到张良，讲述了事情的经过，打算叫张良和他一起离去。\par
\noindent 【7.18.3】他说：“不要跟他们一起死掉。”\par

\begin{itemize}
  \item 时间信息原文：\answerline{7.5cm}（无则填“无”）
  \item 是否为下沉句：\choicebox{是}\quad \choicebox{否}
  \item 是否为插叙/倒叙：\choicebox{是}\quad \choicebox{否}
\end{itemize}

\paragraph{A-1-08 \quad 【7.17.2】}
\textbf{上下文（前五句 / 目标句 / 后五句）：}

\noindent 【7.16.4】秦军官兵多在私下议论说：“章将军等欺骗我们投降诸侯军。如今能够入关破秦，很好；如果不能，诸侯军俘虏我们东去；秦势必把我们的父母妻子全部处死。”\par
\noindent 【7.16.5】诸侯军的将领们暗中听到了他们的打算，报告了项羽。\par
\noindent 【7.16.6】项羽就找来黥布、蒲将军商量说：“秦军官兵还很多，他们心里不服，到了关中不听从命令，事情必然岌岌可危，不如杀掉他们，而只与章邯、长史司马欣、都尉董翳一起入秦。”\par
\noindent 【7.16.7】于是楚军夜间把秦军士卒二十多万人处死掩埋在新安城南。\par
\noindent 【7.17.1】项羽将要攻取秦关中地带。\par
\noindent \textcolor{blue}{\textbf{【7.17.2】函谷关有兵把守，不能进去。}}\par
\noindent 【7.17.3】又听说沛公已经攻破咸阳，项羽大怒，派当阳君等扣关。\par
\noindent 【7.17.4】项羽便进入了函谷关，到达戏水西岸。\par
\noindent 【7.17.5】沛公驻军霸上，没有能够和项羽相见。\par
\noindent 【7.17.6】沛公左司马曹无伤派人对项羽说：“沛公想称王关中，使子婴为相，占有了全部珍宝。”\par
\noindent 【7.17.7】项羽怒气冲天他说：“明天早晨饱餐士卒，将击溃沛公的军队！”\par

\begin{itemize}
  \item 时间信息原文：\answerline{7.5cm}（无则填“无”）
  \item 是否为下沉句：\choicebox{是}\quad \choicebox{否}
  \item 是否为插叙/倒叙：\choicebox{是}\quad \choicebox{否}
\end{itemize}


\newpage
\section{Block B：TimeBlock 推理（Agent 5--6）}
\blocktime
覆盖 Agent 5 时间标志物转换/补全，以及 Agent 6 TimeBlock 前后顺序判断。

\paragraph{B5-01 \quad Agent 5 / 时间标志物补全}
背景信息：前文交代沛公已在汉元年十月先于诸侯到达霸上，秦王投降后，沛公进入咸阳、封闭秦宫府库并回军霸上；随后他采纳建议，派兵把守函谷关，防止诸侯军入关。原文此处只写“十一月间”，应补全为哪一项？


\textbf{材料：}【8.23.3】十一月间，项羽果然率领诸侯军西进，想要入关，而关门闭着。


\begin{enumerate}[label=\Alph*.]
  \item 汉元年十月
  \item 汉元年十一月
  \item 汉元年十二月
  \item 秦二世三年十一月
\end{enumerate}
答案：\answerline{2.2cm}

\paragraph{B5-02 \quad Agent 5 / 时间标志物补全}
背景信息：前文交代项羽在正月自立为西楚霸王，背弃旧约，改立沛公为汉王并令其建都南郑；之后汉王在八月采用韩信计策，从故道回军袭击雍王章邯，并平定雍地。原文此处只写“二年”，应补全为哪一项？


\textbf{材料：}【8.30.1】二年，汉王东出略取城邑，塞王司马欣、翟王董翳、河南王申阳都投降了。


\begin{enumerate}[label=\Alph*.]
  \item 秦二世二年
  \item 汉二年
  \item 汉三年
  \item 汉元年
\end{enumerate}
答案：\answerline{2.2cm}

\paragraph{B6-01 \quad Agent 6 / TimeBlock 顺序判断}
判断 A、B 两个 TimeBlock 的时间先后。


\textbf{材料 A：}【8.16.1】秦二世三年，楚怀王看到项梁的军队被打垮了，心里恐惧，迁离盱台，建都彭城，合并吕臣、项羽的军队，亲自统率。 【8.16.2】以沛公任砀郡长，封为武安侯，统领砀郡的军队。 【8.16.3】封项羽为长安侯，号为鲁公。 【8.16.4】吕臣任司徒，他的父亲吕青作令尹。 【8.17.1】赵多次请求救援，楚怀王就以宋义为上将军，项羽为次将，范增为未将，北上救赵。 【8.17.2】命令沛公西出略地，打入关中。 ……

\textbf{材料 B：}【8.9.1】秦二世元年秋天，陈胜等在蕲县起义，到了陈县自立为王，号称“张楚”。 【8.9.2】各郡县都大多杀死长官，响应陈胜。 【8.9.3】沛县县令恐惧，想要以沛具响应陈胜，主吏萧何、狱掾曹参对他说：“您身为秦朝的官吏，如今要叛秦起事，率领沛县子弟，恐怕他们不愿听命。希望您召集逃亡在外面的人,可以得到几百人。利用这股力量胁持群众，群众不敢不听您的命令。” 【8.9.4】县令就派樊哙去召唤刘季，刘季的队伍已经近百人了。 【8.10.1】于是樊哙跟着刘季来到沛县。 【8.10.2】沛县县令又后悔了，恐怕刘季发生变故，就关闭城门，派人防守，打算杀掉萧何、曹参。 ……


\begin{enumerate}[label=\Alph*.]
  \item A 早于 B
  \item B 早于 A
  \item A 与 B 为同一事件或大致同时
  \item 仅凭材料无法判断
\end{enumerate}
答案：\answerline{2.2cm}

\paragraph{B6-02 \quad Agent 6 / TimeBlock 顺序判断}
判断 A、B 两个 TimeBlock 的时间先后。


\textbf{材料 A：}【8.22.1】沛公的军队先于各路诸侯到达霸上。 【8.22.2】秦王于婴素车白马，用丝带系着脖子，封了皇帝的印玺和符节，在轵道旁投降。 【8.22.3】将领们有的主张杀死秦王。 【8.22.4】沛公说：“当初楚怀王派遣我，本来是因为我能宽大容人。况且人家已经降服，又杀死人家，不吉利。” 【8.22.5】于是就把秦王交给了官吏，向西进入咸阳。 【8.22.6】沛公想要留在宫殿中休息，樊哙、张良劝说后，才封闭了秦宫的贵重珍宝、财物和库房，回军霸上。 ……

\textbf{材料 B：}【8.23.3】项羽果然率领诸侯军西进，想要入关，而关门闭着。


\begin{enumerate}[label=\Alph*.]
  \item A 早于 B
  \item B 早于 A
  \item A 与 B 为同一事件或大致同时
  \item 仅凭材料无法判断
\end{enumerate}
答案：\answerline{2.2cm}


\newpage
\section{Block C：跨文本对齐（Agent 9）}
\blocktime
请选择最可能描述同一历史过程的 source。每题只有一个最佳答案。

\paragraph{C9-01 \quad Agent 9 / 跨文本对齐}
哪一个候选 source 最可能与 target 描述同一历史过程，从而可以向 target 传播时间信息？

\textbf{Target：}\textcolor{blue}{\textbf{【7.3.15】项羽派人搜罗下属各县丁壮，得到精兵八千人。}}

\begin{enumerate}[label=\Alph*.]
  \item
\noindent 【8.14.4】沛公和项羽正在攻打陈留，听说项梁死了，带兵和吕将军一起向东进发。\par
\noindent 【8.14.5】吕臣驻扎在彭城东面，项羽驻扎在彭城西面，沛公驻扎在肠。\par
\noindent 【8.15.1】章邯已经打垮了项梁的军队，以为楚地的敌人不用担心了，就渡过黄河，北进攻打赵地，大破赵军。\par
\noindent 【8.15.2】这个时候，赵歇为赵王，秦将王离围困赵歇于巨鹿城。\par
\noindent 【8.15.3】（被围在巨鹿的军队，）这就是所谓的“河北之军”。\par
\noindent 【8.16.1】秦二世三年，楚怀王看到项梁的军队被打垮了，心里恐惧，迁离盱台，建都彭城，合并吕臣、项羽的军队，亲自统率。\par
\noindent 【8.16.2】以沛公任砀郡长，封为武安侯，统领砀郡的军队。\par
\noindent 【8.16.3】封项羽为长安侯，号为鲁公。\par
\noindent 【8.16.4】吕臣任司徒，他的父亲吕青作令尹。\par
\noindent 【8.17.1】赵多次请求救援，楚怀王就以宋义为上将军，项羽为次将，范增为未将，北上救赵。\par
\noindent 【8.17.2】命令沛公西出略地，打入关中。\par
  \item
\noindent 【8.10.11】在沛县衙门的庭院里祭祀黄帝和蚩尤，又用牲血衅鼓旗。\par
\noindent 【8.10.12】旗子一律红色，因为刘季所杀蛇是白帝的儿子，杀蛇的是赤帝的儿子，所以崇尚赤色。\par
\noindent 【8.10.13】于是少年子弟和有势的官吏，如萧何、曹参、樊哙等人，都为沛公征集兵员，集合了两三千人，攻打胡陵、方与，回军固守丰邑。\par
\noindent 【8.11.1】秦二世二年，陈胜将领周章的军队西至戏水而还。\par
\noindent 【8.11.2】燕、赵、齐、魏都自立为王。\par
\noindent 【8.11.3】项梁、项羽起兵于吴。\par
\noindent 【8.11.4】秦泗水郡郡监平率兵围丰，两天后，沛公出兵应战，打败了秦军。\par
\noindent 【8.11.5】沛公命令雍齿守卫丰邑，自己引兵赴薛。\par
\noindent 【8.11.6】泗水郡郡守壮在薛战败，逃到戚。\par
\noindent 【8.11.7】沛公左司马擒获泗水郡郡守壮，杀死了他。\par
\noindent 【8.11.8】沛公回军亢父，到了方与，没有交战。\par
  \item
\noindent 【8.22.8】沛公派人与秦朝官吏巡行县城乡间，告谕百姓。\par
\noindent 【8.22.9】秦地的百姓大为高兴，争先恐后地拿出牛羊酒食款待士兵。\par
\noindent 【8.22.10】沛公又谦让不肯接受，说：“仓库的谷子很多，不缺乏，不愿破费百姓。”百姓更加高兴，唯恐沛公不做秦王。\par
\noindent 【8.23.1】有人劝沛公说：“秦地比天下富足十倍，地势好。如今听说章邯投降了项羽，项羽就给了雍王的封号，称王于关中。现在即将来到关中就国，你沛公恐怕不能占有这个地方了。应赶快派兵把守函谷关，不让诸侯军进来，逐渐征集关中兵，以加强实力，抵抗诸侯兵。”\par
\noindent 【8.23.2】沛公赞成他的计策，照着做了。\par
\noindent 【8.23.3】十一月间，项羽果然率领诸侯军西进，想要入关，而关门闭着。\par
\noindent 【8.23.4】听说沛公已经平定关中，大怒，派黥布等攻破了函谷关。\par
\noindent 【8.23.5】十二月间，就到了戏水。\par
\noindent 【8.23.6】沛公左司马曹无伤听说项王发怒，\par
\noindent 【8.23.7】要攻打沛公，派人告诉项羽说：“沛公想要称王关中，令子婴为相，珍宝被他全部占有了。”\par
\noindent 【8.23.8】打算以此求得封赏。\par
  \item
\noindent 【8.28.7】齐借兵给陈余，击败了常山王张耳，张耳逃跑归附了汉王。\par
\noindent 【8.28.8】陈余从代接回赵王歇，又立为赵王，赵王就封陈余为代王。\par
\noindent 【8.28.9】项羽大怒，出兵北向击齐。\par
\noindent 【8.29.1】八月，汉王用韩信的计策，从故道回军，袭击雍王章邯。\par
\noindent 【8.29.2】章邯在陈仓迎击汉军，雍王兵败退走，在好畤停下来接战，又失败了，逃到废丘。\par
\noindent 【8.29.3】汉王随即平定了雍地。\par
\noindent 【8.29.4】向东到达咸阳，率军围困雍王于废丘，而派遣将领攻占了陇西、北地、上郡。\par
\noindent 【8.29.5】派将军薛欧、王吸出武关，借助王陵驻扎在南阳的兵力，迎接太公、吕后于沛县。\par
\noindent 【8.29.6】楚听到这一消息，出兵在阳夏阻挡，汉军不能前进。\par
\noindent 【8.29.7】楚让原吴县县令郑昌为韩王，抵抗汉军。\par
\noindent 【8.30.1】二年，汉王东出略取城邑，塞王司马欣、翟王董翳、河南王申阳都投降了。\par
\end{enumerate}
答案：\answerline{2.2cm}

\paragraph{C9-02 \quad Agent 9 / 跨文本对齐}
哪一个候选 source 最可能与 target 描述同一历史过程，从而可以向 target 传播时间信息？

\textbf{Target：}\textcolor{blue}{\textbf{【7.10.1】楚军在定陶打了败仗，楚怀王很恐惧，从盱台前往彭城，合并了项羽、吕臣的军队亲自统率。}}

\begin{enumerate}[label=\Alph*.]
  \item
\noindent 【8.10.11】在沛县衙门的庭院里祭祀黄帝和蚩尤，又用牲血衅鼓旗。\par
\noindent 【8.10.12】旗子一律红色，因为刘季所杀蛇是白帝的儿子，杀蛇的是赤帝的儿子，所以崇尚赤色。\par
\noindent 【8.10.13】于是少年子弟和有势的官吏，如萧何、曹参、樊哙等人，都为沛公征集兵员，集合了两三千人，攻打胡陵、方与，回军固守丰邑。\par
\noindent 【8.11.1】秦二世二年，陈胜将领周章的军队西至戏水而还。\par
\noindent 【8.11.2】燕、赵、齐、魏都自立为王。\par
\noindent 【8.11.3】项梁、项羽起兵于吴。\par
\noindent 【8.11.4】秦泗水郡郡监平率兵围丰，两天后，沛公出兵应战，打败了秦军。\par
\noindent 【8.11.5】沛公命令雍齿守卫丰邑，自己引兵赴薛。\par
\noindent 【8.11.6】泗水郡郡守壮在薛战败，逃到戚。\par
\noindent 【8.11.7】沛公左司马擒获泗水郡郡守壮，杀死了他。\par
\noindent 【8.11.8】沛公回军亢父，到了方与，没有交战。\par
  \item
\noindent 【8.22.8】沛公派人与秦朝官吏巡行县城乡间，告谕百姓。\par
\noindent 【8.22.9】秦地的百姓大为高兴，争先恐后地拿出牛羊酒食款待士兵。\par
\noindent 【8.22.10】沛公又谦让不肯接受，说：“仓库的谷子很多，不缺乏，不愿破费百姓。”百姓更加高兴，唯恐沛公不做秦王。\par
\noindent 【8.23.1】有人劝沛公说：“秦地比天下富足十倍，地势好。如今听说章邯投降了项羽，项羽就给了雍王的封号，称王于关中。现在即将来到关中就国，你沛公恐怕不能占有这个地方了。应赶快派兵把守函谷关，不让诸侯军进来，逐渐征集关中兵，以加强实力，抵抗诸侯兵。”\par
\noindent 【8.23.2】沛公赞成他的计策，照着做了。\par
\noindent 【8.23.3】十一月间，项羽果然率领诸侯军西进，想要入关，而关门闭着。\par
\noindent 【8.23.4】听说沛公已经平定关中，大怒，派黥布等攻破了函谷关。\par
\noindent 【8.23.5】十二月间，就到了戏水。\par
\noindent 【8.23.6】沛公左司马曹无伤听说项王发怒，\par
\noindent 【8.23.7】要攻打沛公，派人告诉项羽说：“沛公想要称王关中，令子婴为相，珍宝被他全部占有了。”\par
\noindent 【8.23.8】打算以此求得封赏。\par
  \item
\noindent 【8.14.4】沛公和项羽正在攻打陈留，听说项梁死了，带兵和吕将军一起向东进发。\par
\noindent 【8.14.5】吕臣驻扎在彭城东面，项羽驻扎在彭城西面，沛公驻扎在肠。\par
\noindent 【8.15.1】章邯已经打垮了项梁的军队，以为楚地的敌人不用担心了，就渡过黄河，北进攻打赵地，大破赵军。\par
\noindent 【8.15.2】这个时候，赵歇为赵王，秦将王离围困赵歇于巨鹿城。\par
\noindent 【8.15.3】（被围在巨鹿的军队，）这就是所谓的“河北之军”。\par
\noindent 【8.16.1】秦二世三年，楚怀王看到项梁的军队被打垮了，心里恐惧，迁离盱台，建都彭城，合并吕臣、项羽的军队，亲自统率。\par
\noindent 【8.16.2】以沛公任砀郡长，封为武安侯，统领砀郡的军队。\par
\noindent 【8.16.3】封项羽为长安侯，号为鲁公。\par
\noindent 【8.16.4】吕臣任司徒，他的父亲吕青作令尹。\par
\noindent 【8.17.1】赵多次请求救援，楚怀王就以宋义为上将军，项羽为次将，范增为未将，北上救赵。\par
\noindent 【8.17.2】命令沛公西出略地，打入关中。\par
  \item
\noindent 【8.28.7】齐借兵给陈余，击败了常山王张耳，张耳逃跑归附了汉王。\par
\noindent 【8.28.8】陈余从代接回赵王歇，又立为赵王，赵王就封陈余为代王。\par
\noindent 【8.28.9】项羽大怒，出兵北向击齐。\par
\noindent 【8.29.1】八月，汉王用韩信的计策，从故道回军，袭击雍王章邯。\par
\noindent 【8.29.2】章邯在陈仓迎击汉军，雍王兵败退走，在好畤停下来接战，又失败了，逃到废丘。\par
\noindent 【8.29.3】汉王随即平定了雍地。\par
\noindent 【8.29.4】向东到达咸阳，率军围困雍王于废丘，而派遣将领攻占了陇西、北地、上郡。\par
\noindent 【8.29.5】派将军薛欧、王吸出武关，借助王陵驻扎在南阳的兵力，迎接太公、吕后于沛县。\par
\noindent 【8.29.6】楚听到这一消息，出兵在阳夏阻挡，汉军不能前进。\par
\noindent 【8.29.7】楚让原吴县县令郑昌为韩王，抵抗汉军。\par
\noindent 【8.30.1】二年，汉王东出略取城邑，塞王司马欣、翟王董翳、河南王申阳都投降了。\par
\end{enumerate}
答案：\answerline{2.2cm}

\paragraph{C9-03 \quad Agent 9 / 跨文本对齐}
哪一个候选 source 最可能与 target 描述同一历史过程，从而可以向 target 传播时间信息？

\textbf{Target：}\textcolor{blue}{\textbf{【7.17.2】函谷关有兵把守，不能进去。}}

\begin{enumerate}[label=\Alph*.]
  \item
\noindent 【8.22.8】沛公派人与秦朝官吏巡行县城乡间，告谕百姓。\par
\noindent 【8.22.9】秦地的百姓大为高兴，争先恐后地拿出牛羊酒食款待士兵。\par
\noindent 【8.22.10】沛公又谦让不肯接受，说：“仓库的谷子很多，不缺乏，不愿破费百姓。”百姓更加高兴，唯恐沛公不做秦王。\par
\noindent 【8.23.1】有人劝沛公说：“秦地比天下富足十倍，地势好。如今听说章邯投降了项羽，项羽就给了雍王的封号，称王于关中。现在即将来到关中就国，你沛公恐怕不能占有这个地方了。应赶快派兵把守函谷关，不让诸侯军进来，逐渐征集关中兵，以加强实力，抵抗诸侯兵。”\par
\noindent 【8.23.2】沛公赞成他的计策，照着做了。\par
\noindent 【8.23.3】十一月间，项羽果然率领诸侯军西进，想要入关，而关门闭着。\par
\noindent 【8.23.4】听说沛公已经平定关中，大怒，派黥布等攻破了函谷关。\par
\noindent 【8.23.5】十二月间，就到了戏水。\par
\noindent 【8.23.6】沛公左司马曹无伤听说项王发怒，\par
\noindent 【8.23.7】要攻打沛公，派人告诉项羽说：“沛公想要称王关中，令子婴为相，珍宝被他全部占有了。”\par
\noindent 【8.23.8】打算以此求得封赏。\par
  \item
\noindent 【8.14.4】沛公和项羽正在攻打陈留，听说项梁死了，带兵和吕将军一起向东进发。\par
\noindent 【8.14.5】吕臣驻扎在彭城东面，项羽驻扎在彭城西面，沛公驻扎在肠。\par
\noindent 【8.15.1】章邯已经打垮了项梁的军队，以为楚地的敌人不用担心了，就渡过黄河，北进攻打赵地，大破赵军。\par
\noindent 【8.15.2】这个时候，赵歇为赵王，秦将王离围困赵歇于巨鹿城。\par
\noindent 【8.15.3】（被围在巨鹿的军队，）这就是所谓的“河北之军”。\par
\noindent 【8.16.1】秦二世三年，楚怀王看到项梁的军队被打垮了，心里恐惧，迁离盱台，建都彭城，合并吕臣、项羽的军队，亲自统率。\par
\noindent 【8.16.2】以沛公任砀郡长，封为武安侯，统领砀郡的军队。\par
\noindent 【8.16.3】封项羽为长安侯，号为鲁公。\par
\noindent 【8.16.4】吕臣任司徒，他的父亲吕青作令尹。\par
\noindent 【8.17.1】赵多次请求救援，楚怀王就以宋义为上将军，项羽为次将，范增为未将，北上救赵。\par
\noindent 【8.17.2】命令沛公西出略地，打入关中。\par
  \item
\noindent 【8.21.7】乘胜追击，彻底打垮了秦军。\par
\noindent 【8.22.1】汉元年十月，沛公的军队先于各路诸侯到达霸上。\par
\noindent 【8.22.2】秦王于婴素车白马，用丝带系着脖子，封了皇帝的印玺和符节，在轵道旁投降。\par
\noindent 【8.22.3】将领们有的主张杀死秦王。\par
\noindent 【8.22.4】沛公说：“当初楚怀王派遣我，本来是因为我能宽大容人。况且人家已经降服，又杀死人家，不吉利。”\par
\noindent 【8.22.5】于是就把秦王交给了官吏，向西进入咸阳。\par
\noindent 【8.22.6】沛公想要留在宫殿中休息，樊哙、张良劝说后，才封闭了秦宫的贵重珍宝、财物和库房，回军霸上。\par
\noindent 【8.22.7】召集各县的父老、豪杰说：“父老们苦于秦朝的严刑峻法已经很久了，诽谤朝政的要灭族，相聚议论的要在街市上处斩。我和诸侯们约定，先入关的在关中称王，我应当称王关中。同父老们约定，法律只有三章：杀人的处死，伤人和抢劫的处以与所犯罪相当的刑罚。其余的秦朝法律全部废除。官吏和百姓都要安居如故。我所以到这里来，是为父老们除害，不会有欺凌暴虐的行为，不要害怕。我所以回军霸上，是等待诸侯们到来制定共同遵守的纪律。”\par
\noindent 【8.22.8】沛公派人与秦朝官吏巡行县城乡间，告谕百姓。\par
\noindent 【8.22.9】秦地的百姓大为高兴，争先恐后地拿出牛羊酒食款待士兵。\par
\noindent 【8.22.10】沛公又谦让不肯接受，说：“仓库的谷子很多，不缺乏，不愿破费百姓。”百姓更加高兴，唯恐沛公不做秦王。\par
  \item
\noindent 【8.28.7】齐借兵给陈余，击败了常山王张耳，张耳逃跑归附了汉王。\par
\noindent 【8.28.8】陈余从代接回赵王歇，又立为赵王，赵王就封陈余为代王。\par
\noindent 【8.28.9】项羽大怒，出兵北向击齐。\par
\noindent 【8.29.1】八月，汉王用韩信的计策，从故道回军，袭击雍王章邯。\par
\noindent 【8.29.2】章邯在陈仓迎击汉军，雍王兵败退走，在好畤停下来接战，又失败了，逃到废丘。\par
\noindent 【8.29.3】汉王随即平定了雍地。\par
\noindent 【8.29.4】向东到达咸阳，率军围困雍王于废丘，而派遣将领攻占了陇西、北地、上郡。\par
\noindent 【8.29.5】派将军薛欧、王吸出武关，借助王陵驻扎在南阳的兵力，迎接太公、吕后于沛县。\par
\noindent 【8.29.6】楚听到这一消息，出兵在阳夏阻挡，汉军不能前进。\par
\noindent 【8.29.7】楚让原吴县县令郑昌为韩王，抵抗汉军。\par
\noindent 【8.30.1】二年，汉王东出略取城邑，塞王司马欣、翟王董翳、河南王申阳都投降了。\par
\end{enumerate}
答案：\answerline{2.2cm}

\paragraph{C9-04 \quad Agent 9 / 跨文本对齐}
哪一个候选 source 最可能与 target 描述同一历史过程，从而可以向 target 传播时间信息？

\textbf{Target：}\textcolor{blue}{\textbf{【7.27.1】这时，汉王回军平定了三秦。}}

\begin{enumerate}[label=\Alph*.]
  \item
\noindent 【8.22.8】沛公派人与秦朝官吏巡行县城乡间，告谕百姓。\par
\noindent 【8.22.9】秦地的百姓大为高兴，争先恐后地拿出牛羊酒食款待士兵。\par
\noindent 【8.22.10】沛公又谦让不肯接受，说：“仓库的谷子很多，不缺乏，不愿破费百姓。”百姓更加高兴，唯恐沛公不做秦王。\par
\noindent 【8.23.1】有人劝沛公说：“秦地比天下富足十倍，地势好。如今听说章邯投降了项羽，项羽就给了雍王的封号，称王于关中。现在即将来到关中就国，你沛公恐怕不能占有这个地方了。应赶快派兵把守函谷关，不让诸侯军进来，逐渐征集关中兵，以加强实力，抵抗诸侯兵。”\par
\noindent 【8.23.2】沛公赞成他的计策，照着做了。\par
\noindent 【8.23.3】十一月间，项羽果然率领诸侯军西进，想要入关，而关门闭着。\par
\noindent 【8.23.4】听说沛公已经平定关中，大怒，派黥布等攻破了函谷关。\par
\noindent 【8.23.5】十二月间，就到了戏水。\par
\noindent 【8.23.6】沛公左司马曹无伤听说项王发怒，\par
\noindent 【8.23.7】要攻打沛公，派人告诉项羽说：“沛公想要称王关中，令子婴为相，珍宝被他全部占有了。”\par
\noindent 【8.23.8】打算以此求得封赏。\par
  \item
\noindent 【8.25.2】楚怀王说：“按照原来的约定办。”\par
\noindent 【8.25.3】项羽怨恨楚怀王不肯让他与沛公一起西进入关，而派他北上救赵，\par
\noindent 【8.25.4】在天下诸侯争夺称王关中的约定中落在后面，他就说：“怀王这个人，我家项梁所立，没有什么功劳，凭什么主持约定。本来安定天下的，是诸位将领和我项籍。”\par
\noindent 【8.25.5】就假意推尊楚怀王为义帝，实际上不听从他的命令。\par
\noindent 【8.26.1】正月，项羽自立为西楚霸王，在梁、楚地区的九个郡称王，建都彭城。\par
\noindent 【8.26.2】背弃原来的约定，改立沛公为汉王，在巴、蜀、汉中称玉，建都南郑。\par
\noindent 【8.26.3】把关中瓜分为三，封立秦朝的三个将领：章邯为雍王，建都废丘：司马欣为塞王，建都栎阳；董翳为翟王，建都高奴。\par
\noindent 【8.26.4】封楚将瑕丘申阳为河南王，建都洛阳。\par
\noindent 【8.26.5】封赵将司马印为殷王，建都朝歌。\par
\noindent 【8.26.6】赵王歇迁徙代地称王。\par
\noindent 【8.26.7】封赵将张耳为常山王，建都襄国。\par
  \item
\noindent 【8.14.4】沛公和项羽正在攻打陈留，听说项梁死了，带兵和吕将军一起向东进发。\par
\noindent 【8.14.5】吕臣驻扎在彭城东面，项羽驻扎在彭城西面，沛公驻扎在肠。\par
\noindent 【8.15.1】章邯已经打垮了项梁的军队，以为楚地的敌人不用担心了，就渡过黄河，北进攻打赵地，大破赵军。\par
\noindent 【8.15.2】这个时候，赵歇为赵王，秦将王离围困赵歇于巨鹿城。\par
\noindent 【8.15.3】（被围在巨鹿的军队，）这就是所谓的“河北之军”。\par
\noindent 【8.16.1】秦二世三年，楚怀王看到项梁的军队被打垮了，心里恐惧，迁离盱台，建都彭城，合并吕臣、项羽的军队，亲自统率。\par
\noindent 【8.16.2】以沛公任砀郡长，封为武安侯，统领砀郡的军队。\par
\noindent 【8.16.3】封项羽为长安侯，号为鲁公。\par
\noindent 【8.16.4】吕臣任司徒，他的父亲吕青作令尹。\par
\noindent 【8.17.1】赵多次请求救援，楚怀王就以宋义为上将军，项羽为次将，范增为未将，北上救赵。\par
\noindent 【8.17.2】命令沛公西出略地，打入关中。\par
  \item
\noindent 【8.28.7】齐借兵给陈余，击败了常山王张耳，张耳逃跑归附了汉王。\par
\noindent 【8.28.8】陈余从代接回赵王歇，又立为赵王，赵王就封陈余为代王。\par
\noindent 【8.28.9】项羽大怒，出兵北向击齐。\par
\noindent 【8.29.1】八月，汉王用韩信的计策，从故道回军，袭击雍王章邯。\par
\noindent 【8.29.2】章邯在陈仓迎击汉军，雍王兵败退走，在好畤停下来接战，又失败了，逃到废丘。\par
\noindent 【8.29.3】汉王随即平定了雍地。\par
\noindent 【8.29.4】向东到达咸阳，率军围困雍王于废丘，而派遣将领攻占了陇西、北地、上郡。\par
\noindent 【8.29.5】派将军薛欧、王吸出武关，借助王陵驻扎在南阳的兵力，迎接太公、吕后于沛县。\par
\noindent 【8.29.6】楚听到这一消息，出兵在阳夏阻挡，汉军不能前进。\par
\noindent 【8.29.7】楚让原吴县县令郑昌为韩王，抵抗汉军。\par
\noindent 【8.30.1】二年，汉王东出略取城邑，塞王司马欣、翟王董翳、河南王申阳都投降了。\par
\end{enumerate}
答案：\answerline{2.2cm}


\end{document}
